\clearpage
\section{Appendix: Liturgical Resolution}
\label{sec:liturgical-resolution}

{\small
\textbf{Celebration.} The First Sunday of Advent, a Sunday of Advent at
place~2 of the table of liturgical days; violet; no Gloria; Creed. The formula
is \bid{PC-S01-B}, the Year~B target of the parent \bid{PC-S01}; it has a
single form, and every reading of Lectionary no.~2 has one form only. Advent
begins with Evening Prayer~I of this Sunday, so a Mass anticipated on Saturday
evening, 28~November, uses the same texts and readings.

\textbf{Books, calendar and territory.} The \work{Roman Missal}, Third
Edition, for use in the dioceses of the United States of America (2011), whose
Latin base is the \work{Missale Romanum}, editio typica tertia (2002), as
reprinted with emendations in 2008; the \work{Lectionary for Mass for Use in
the Dioceses of the United States of America}, second typical edition,
volume~I, no.~2, which follows the \work{Ordo lectionum Missae}, editio typica
altera (1981); the \work{General Instruction of the Roman Missal} in the
United States edition of 2011 with the emendations of 2021; the
\work{Universal Norms on the Liturgical Year and the Calendar} and the General
Roman Calendar, in the Latin printed in the 2002 Missal; and the
\work{Liturgical Calendar for the Dioceses of the United States of America} for
2026 and for 2027. Only the national calendar is applied. No diocesan,
religious, parish, titular, dedication or patronal calendar was supplied, and
none is assumed.

\textbf{Cycle and occurrence.} Under \work{Universal Norms} 40 Advent begins
on the Sunday that falls on or nearest 30~November; in 2026 that is Sunday
29~November. The Sunday opens the liturgical year 2027 and Year~B of the
Sunday cycle, which runs to 21~November 2027; the national calendar gives
Lectionary no.~2 and Psalter Week~I. The Sundays of Advent take precedence
over every feast of the Lord and every solemnity (\work{Universal Norms} 5),
and the General Roman Calendar inscribes nothing on 29~November. The Feast of
Saint Andrew, Apostle, falls on Monday 30~November and contributes nothing
to this Mass.

\textbf{The Missal's texts and the three cycles.} The antiphons, prayers,
Preface and blessing belong to the First Sunday of Advent and are common to
Years A, B and~C; their text is recorded once, in the collection's record for
that Sunday. The readings, psalm and Alleluia verse are those of Year~B. In
Advent the readings are chosen for the season and are not read in course
(\work{Praenotanda} nos.~67 and~93). The Prayer over the Offerings and the
Prayer after Communion are used again on Monday of the First Week of Advent.

\textbf{Branches.} No selection among the following was supplied, and none is
inferred. A claim resting on the wording of an antiphon, the psalm or the
Alleluia verse does not hold where another approved chant replaces it.
\par}

\begin{branchtable}
\bid{entrance-or-communion-chant-substitution} & \work{General Instruction} 48
and 87: the Missal or \work{Graduale Romanum} antiphon, the \work{Graduale
Simplex}, a chant from another collection of psalms and antiphons approved by
the Conference of Bishops or the Diocesan Bishop, or another suitable
liturgical chant likewise approved & permitted & Entrance; Communion &
Unselected; the Missal antiphons are the texts studied.\\
\bid{psalm-or-acclamation-chant-substitution} & \work{General Instruction} 61
and 62: when the psalm is sung, a common response and psalm for the season
(for Advent the \work{Ordo}'s common response is \latin{Veni, Domine, ad
liberandum nos}, with Ps 25 (24) or Ps 85 (84)); the \work{Graduale} forms or
an antiphon and psalm from another approved collection, never a song or hymn
in place of the psalm & permitted & Responsorial Psalm; Alleluia verse &
Unselected; the Lectionary's psalm and verse are the texts studied.\\
\bid{offertory-chant} & \work{General Instruction} 74 & permitted &
Preparation of the gifts & Unselected; the Missal prints no Offertory
antiphon.\\
\bid{sprinkling-rite} & \work{General Instruction} 51: on Sundays the blessing
and sprinkling of water may replace the Penitential Act & ritual substitution
& Introductory Rites & Unselected; it brings its own texts.\\
\bid{eucharistic-prayer-with-appointed-preface} & The Missal's rubric
appoints Preface~I of Advent; \work{General Instruction} 364--365 a--c &
permitted & Preface; Eucharistic Prayer & Unselected. Prayer I, II or III may
be said with the appointed Preface. If Prayer~II is said with its own Preface,
Preface~I of Advent is not said.\\
\bid{eucharistic-prayer-iv-with-preface} & \work{General Instruction} 365~d &
not admitted & Eucharistic Prayer & Prayer~IV, with its invariable Preface,
is for Masses without a proper Preface and for Sundays in Ordinary Time;
neither condition holds.\\
\bid{reconciliation-or-various-needs-eucharistic-prayer} & The Missal's
rubrics for the Eucharistic Prayers for Reconciliation and for Various Needs
& not applicable & Eucharistic Prayer & Neither rubric names a Sunday of Advent
or an Advent Preface.\\
\bid{concluding-blessing-form} & The Missal's rubric for this Mass: the solemn
blessing for Advent may be used; the solemn blessings and prayers over the
people are at the priest's discretion & permitted & Concluding Rites &
Unselected. Where the solemn blessing is not used, the simple blessing of the
Order of Mass is given.\\
\bid{advent-wreath-blessing} & \work{Book of Blessings} nos. 1509--1520, read
in a published transcription: on the First Sunday of Advent the wreath may be
blessed within Mass, the prayer of blessing concluding the general
intercessions & permitted & Universal Prayer & Unselected. It adds a blessing
and replaces no text of this Mass.\\
\bid{ritual-or-various-needs-mass} & \work{General Instruction} 372 and 374:
Ritual Masses are prohibited on the Sundays of Advent, and a Mass for a
graver need is excluded on them & not admitted & Whole Mass & Excluded on this date.\\
\bid{local-proper-calendar} & \work{Universal Norms} 5, 59 and 60: a proper
solemnity of a particular calendar ranks below the Sundays of Advent and is
transferred to the Monday & not admitted & Whole Mass & No local celebration
displaces this Sunday. A local solemnity moved to Monday 30~November would
meet the Feast of Saint Andrew, which is a question for a local
resolution.\\
\end{branchtable}
